\chapter{Zasnova sistema}
\label{ch:zasnova}

\section{Funkcionalne in nefunkcionalne zahteve}

Izhodišče zasnove je samostojna namizna naprava, ki zvočno dogajanje pretvori
v vidno opozorilo. Primarni uporabni scenarij je pomoč gluhim in naglušnim
osebam pri zaznavanju izbranih nevarnih zvokov. Naprava zato ne sme zahtevati
telefona, osebnega računalnika ali dostopa do interneta. Zahteve so zbrane v
tabeli \ref{tab:zahteve-sistema}.

\begin{table}[htbp]
  \centering
  \caption{Glavne zahteve sistema.}
  \label{tab:zahteve-sistema}
  \small
  \begin{tabular}{p{0.22\textwidth}p{0.63\textwidth}}
    \toprule
    Vrsta & Zahteva \\
    \midrule
    Funkcionalna & Neprekinjen zajem zvoka z vgrajenega mikrofona. \\
    Funkcionalna & Prikaz pasjega laježa, razbitja stekla, strela, sirene,
                   govora in drugih zvokov. \\
    Funkcionalna & Ločeni stanji za nezanesljivo napoved in neaktiven vhod. \\
    Funkcionalna & Prikaz treh najvišjih ocen in zgodovine potrjenih dogodkov. \\
    Nefunkcionalna & Sprotno izvajanje brez izgube zvočnih vzorcev. \\
    Nefunkcionalna & Delovanje brez omrežne povezave in računalniškega oblaka. \\
    Nefunkcionalna & Brez trajnega shranjevanja surovega zvoka v običajnem načinu. \\
    Nefunkcionalna & Ponovljiv zagon podpisane aplikacije ter sledljiva različica
                     modela in uteži. \\
    \bottomrule
  \end{tabular}
\end{table}

Zadnji dve zahtevi sta povezani z zasebnostjo in zanesljivostjo. Mikrofon mora
zaradi namena sistema delovati neprekinjeno, vendar zvočni vzorci ostanejo v
začasnih medpomnilnikih na napravi in se po obdelavi prepišejo. Poseben
diagnostični način za zajem zvoka je bil uporabljen samo med razvojem in ni
vključen v običajni zagon končnega izdelka.

\section{Strojna platforma STM32N6570-DK}

Razvojna plošča STM32N6570-DK vsebuje mikrokrmilnik STM32N657 z jedrom Arm
Cortex-M55 in namenskim nevronskim pospeševalnikom Neural-ART. Na isti plošči
so vgrajeni digitalni mikrofon IMP34DT05, zunanji bliskovni pomnilnik,
pomnilnik HyperRAM in zaslon ločljivosti $800 \times 480$ slikovnih točk
\cite{st2024stm32n6570dk}. Zaradi te kombinacije za osnovno delovanje prototipa
niso potrebni zunanji mikrofon, grafični krmilnik ali dodaten računalnik.

Procesorsko jedro vodi zajem, pripravo značilk, odločitveni filter in risanje
uporabniškega vmesnika. Neural-ART izvaja računsko najzahtevnejše dele
nevronske mreže, zunanji bliskovni pomnilnik hrani podpisano aplikacijo in
uteži, HyperRAM pa dva zaslonska medpomnilnika. Takšna razdelitev prepreči, da
bi prikaz porabil notranji pomnilnik, potreben za program in začasne aktivacije
modela.

\begin{figure}[htbp]
  \centering
  \includegraphics[width=0.86\textwidth]{20260809_185828.jpg}
  \caption{Razvojna plošča STM32N6570-DK z vidnim mikrokrmilnikom in
  priključki, uporabljena za izvedbo sistema.}
  \label{fig:plosca}
\end{figure}

\section{Arhitektura podatkovnega toka}

Podatkovni tok je prikazan na sliki \ref{fig:podatkovni-tok}. Vsak blok ima
jasno nalogo in podatki med običajnim delovanjem ne zapustijo naprave.

\begin{figure}[htbp]
  \centering
  \fbox{\begin{minipage}{0.90\textwidth}
    \centering\small
    \textbf{Vgrajeni mikrofon}
    $\longrightarrow$
    \textbf{digitalni filter in neposredni dostop do pomnilnika}
    $\longrightarrow$
    \textbf{šestnajstbitni vzorci}
    \\[0.8em]
    $\longrightarrow$
    \textbf{logaritemski melov spektrogram}
    $\longrightarrow$
    \textbf{osembitni vhodni tenzor}
    $\longrightarrow$
    \textbf{YAMNet-1024 na Neural-ART}
    \\[0.8em]
    $\longrightarrow$
    \textbf{časovno filtriranje in zavračanje}
    $\longrightarrow$
    \textbf{zaslon in diagnostični zapis}
  \end{minipage}}
  \caption{Poenostavljen tok podatkov od zvoka v prostoru do lokalnega
  vizualnega opozorila.}
  \label{fig:podatkovni-tok}
\end{figure}

Digitalni mikrofon oddaja pulzno-gostotni tok, ki ga strojni filter pretvori v
zvočne vzorce s frekvenco \SI{16}{\kilo\hertz}. Neposredni dostop do pomnilnika
jih zapisuje brez posredovanja procesorskega jedra pri vsakem vzorcu.
Predobdelava iz skoraj ene sekunde zvoka pripravi 64-pasovni logaritemski
melov spektrogram. Ta se kvantizira in neposredno zapiše v vhodni pomnilnik
omrežja. Sedem ocen modela se nato pretvori v šest uporabniških razredov,
časovni filter pa preprečuje, da bi kratka nihanja takoj postala opozorilo.

\section{Programska zgradba in zagonski tok}

Program temelji na referenčni aplikaciji za obdelavo zvoka podjetja
STMicroelectronics, namenjeni mikrokrmilniku STM32N6 \cite{ststm32n6audio}.
Referenčna osnova zagotavlja začetno nastavitev ure,
mikrofona, pomnilnikov in pospeševalnika, nad njo pa so bili dodani ciljni
model, odločitveni filter, diagnostični zapisi ter lasten prikaz. Končna
različica deluje brez operacijskega sistema; glavna zanka obdela nov zvočni
blok, ko ga je v krožnem medpomnilniku dovolj.

Po ponastavitvi začetni nalagalnik preveri podpis aplikacije v zunanjem
bliskovnem pomnilniku in jo prenese v notranji pomnilnik za izvajanje. Uteži
modela ostanejo v ločeni particiji zunanjega bliskovnega pomnilnika. Aplikacija
nato inicializira zunanji pomnilnik, zaslon, zvočni zajem, predobdelavo in
izvajalno okolje nevronske mreže. Takšna zgradba omogoča ločeno posodobitev
aplikacije in uteži ter nedvoumno preverjanje obeh datotek s kontrolnima
vsotama.

\section{Nabor razredov in politika odločanja}

Nevronska mreža ima sedem učnih izhodov: pasji lajež, razbitje stekla, strel,
drugi zvoki, sirena, govor in nevihta. Končni uporabniški vmesnik prikazuje
šest razredov. Nevihta po končni odločitvi o namenu izdelka ni obravnavana kot
samostojen nevarnostni dogodek; njena ocena se pred glajenjem prišteje razredu
\emph{drugo}. S tem se ohrani že naučeni model, uporabnik pa nevihte ne vidi
kot posebnega opozorila.

Nevarnostni razredi so pasji lajež, razbitje stekla, strel in sirena. Govor in
drugo sta informacijska oziroma varovalna razreda. Njuna naloga je poimenovati
pogoste nenevarne vhode, ki bi sicer lahko tekmovali z nevarnostnimi razredi.
Razred drugo vključuje raznolike zvoke, med njimi ploskanje, aplavz, zvonce,
kratke poke, gospodinjske zvoke, veter in glasbo.

Najvišja surova ocena modela še ni končna odločitev. Ocene se gladijo skozi
čas, vsak razred ima svoj prag in zahtevano število zaporednih potrditev.
\emph{Neznano} pomeni aktiven zvok brez dovolj zanesljivega izhoda, medtem ko
\emph{čakanje} pomeni, da vhod ni presegel energijskega praga. Neznano in
čakanje zato nista dodatna razreda nevronske mreže. Podrobnosti odločitvenega
filtra so opisane v poglavju \ref{ch:implementacija}.
